\label{sec:discussion}

The current results already show that resource-constrained LLM selection is a
multi-objective engineering problem. Qwen3-8B is the strongest exported int4
public benchmark row, but it is not automatically the best deployment choice
for every railway workload. Qwen3-4B uses less memory and remains competitive
on HumanEval and C-Eval, while the 7B-class Qwen2.5 and DeepSeek efficiency
exports provide higher int4 throughput at a larger memory cost. The final
recommendation must therefore combine public quality, railway-domain quality,
latency, and VRAM margin.

The railway-domain results also show why public benchmarks are insufficient.
The exported Qwen3-8B row is clearly stronger than Qwen3-4B on public tasks,
yet both models struggle on Domain-RegQA. Evidence-grounded regulatory QA is
not solved by general benchmark strength alone. This supports the paper's
planned emphasis on domain-specific evaluation and data-composition ablation.

The QLoRA artifacts are promising but incomplete. Adapter-A training works for
Qwen2.5-7B-Instruct and Qwen3-8B within the 24GB constraint. The Qwen2.5
Adapter-A evaluation preserves broad MMLU and C-Eval performance, but its
Domain-QA exact match is still zero and outputs are often too long. This
suggests that full-sequence adaptation on Domain-QA may not be sufficient.
Completion-only masking and mixed regulatory evidence data are plausible next
steps, which is why Adapter-C and Adapter-D are retained as required
placeholders.

Several limitations remain. First, the final ten-model int4 sweep is not yet
complete. Second, the current domain comparison covers only a subset of the
formal model pool. Third, exact match is too strict for some bilingual
translation cases, so the final paper should include qualitative examples and
error categories. Fourth, all measurements target one GPU class and one local
software stack; performance may differ on newer consumer GPUs, datacenter GPUs,
or CPU-only edge machines. Finally, the paper is an application-oriented
evaluation and adaptation study rather than a new compression or training
algorithm.

Despite these limitations, the reorganized draft provides a concrete path to a
complete paper. The remaining work is well bounded: run the missing ten-model
baseline rows, finish the Domain-RegQA and mixed-adapter evaluations, regenerate
the tables from exported artifacts, and replace the placeholders with final
numbers and error analysis.
